\begin{claim}{C.7a}{Proved}
Let $U=\mathbf P^1_{\Q_5}\setminus\{0,1,\infty\}$,
$F=\Q_5(\mu_\infty)$, and $H=G_F$. For either the ordinary point $2$
or the unit tangent $\overrightarrow{01}$, identify the geometric maximal
pro-five group $P$ with the free pro-five group on cusp generators $a,b$
at $0,1$. Put
\[
 N=\ker(P\longrightarrow\Z_5),\qquad a\longmapsto1,\quad b\longmapsto0,
 \qquad Q=P/\overline{[N,N]}.
\]
Then the images of the \emph{full profinite} fixed groups satisfy
\[
 \operatorname{im}\bigl(\Pi^{s_2(H)}\to P\bigr)=1,
 \qquad
 \operatorname{im}\bigl(\Pi^{s_{\overrightarrow{01}}(H)}\to Q\bigr)
 =\overline{\langle a\rangle}.
\]
In particular the ordinary-point full fixed group is contained in
$\ker(\Pi\to\Pi^{(5)})$. The respective images in
$P^{\rm ab}=\Z_5a\oplus\Z_5b$ are zero and the cusp-$0$ line. These are upper bounds obtained using
wild geometric covers; this quotient argument alone does not determine
either full fixed group. The full equalities are proved independently
in C.7h below.
\end{claim}
\begin{proof}
We first give the arithmetic and class-two facts used in the proof.
The extension $F/\Q_5$ is abelian, since it is generated by roots of unity.
If $M\subset F$ is finite over $\Q_5$, restriction is injective on
\[
 H^1(G_M,\Q_5(1))\longrightarrow H^1(H,\Q_5(1)).                 \tag{C.7a.1}
\]
Indeed the inflation--restriction kernel is
$H^1(\operatorname{Gal}(F/M),\Q_5(1))$. Write $\chi$ for the five-adic
cyclotomic character, and choose $\gamma$ in this abelian quotient with
$\chi(\gamma)\ne1$. For a continuous cocycle $c$, commutativity gives
\[
 (\chi(\gamma)-1)c(\sigma)=(\chi(\sigma)-1)c(\gamma).
\]
Thus $c$ is the coboundary of $c(\gamma)/(\chi(\gamma)-1)$.
Local Kummer theory identifies $H^1(G_M,\Q_5(1))$ with
$\widehat{M^\times}^{\,(5)}\otimes_{\Z_5}\Q_5$. On units this identifies
classes with their five-adic logarithms; a nonzero valuation also gives
a nonzero class. These assertions follow directly from the exact Kummer
sequence, Hilbert 90, and the logarithm isomorphism on a sufficiently
small principal-unit subgroup. We write $\kappa(u)$ for the restriction
to $H$ of the Kummer class of $u\in M^\times$.

For a point $t\in M\setminus\{0,1\}$ of the tripod, any $H$-fixed
geometric element with pro-five abelian coordinates $(A,B)$ satisfies
\[
 \kappa(t) B-\kappa(1-t) A=0
 \quad\hbox{in }H^1(H,\Q_5(1)).                              \tag{C.7a.2}
\]
Here the coordinates are the geometric winding numbers at $0$ and $1$.
To verify the formula, use the standard tangential action
\[
 \sigma(a)=a^{\chi(\sigma)},\qquad
 \sigma(b)=f_\sigma^{-1}b^{\chi(\sigma)}f_\sigma,
 \qquad f_\sigma\in\overline{[P,P]}.
\]
This is the explicit formula of
\citet[\S2, equation (2.1) and the following paragraph]{BC_NakamuraSchneps2000},
restricted from a decomposition group over $\Q$ to $G_M$.
On $H$ it acts trivially on $P/\Gamma_3P$.
Changing from the tangent to $t$ conjugates this action by a path cocycle
whose abelianization is $(\kappa(t),\kappa(1-t))$; this also follows by
applying the covers $z^{5^r}=x$ and $z^{5^r}=1-x$ to that path.
In the class-two free group the commutator of elements with coordinates
$(u,v)$ and $(A,B)$ has exponent $uB-vA$ in the central generator $[a,b]$.
Since $\Gamma_2P/\Gamma_3P\simeq\Z_5$, fixedness gives (C.7a.2).
The sign changes arising from a reversed path convention do not change
this zero equation.

We apply (C.7a.2) first at $t=2$. The class $\kappa(-1)$ is zero,
whereas $\kappa(2)\ne0$ by (C.7a.1) and $\log_5(2)\ne0$.
Hence a full $H$-fixed element $g$ has $B=0$.
Now consider the degree-two finite etale cover
\[
 U_2=\mathbf P^1\setminus\{0,1,-1,\infty\}
       \longrightarrow U,\qquad x=z^2.
\]
A chosen lift $\alpha=\sqrt2$ of $2$ lies in the unramified quadratic
extension $M/\Q_5$, and therefore in $F$: prime-to-five roots of unity
generate the maximal unramified extension. The element $g^2$ belongs to
the geometric subgroup of this cover and is fixed by the lifted section
at $\alpha$. Forget the cusp $-1$, obtaining the tripod with ordinary
base point $\alpha$.
The two classes $\kappa(\alpha)$ and $\kappa(1-\alpha)$ are linearly
independent over $\Q_5$ on $H$. In fact $\log_5(\alpha)=\frac12\log_5(2)$
is a nonzero element of $\Q_5$, whereas
\[
 \operatorname{Tr}_{M/\Q_5}\log_5(1-\alpha)
 =\log_5\bigl((1-\alpha)(1+\alpha)\bigr)=\log_5(-1)=0.
\]
The latter logarithm is nonzero: $1-\sqrt2$ is a unit and is not a root
of unity, as is seen from either of its real algebraic conjugates.
Trace on $\Q_5$ is multiplication by $2$, proving independence before
restriction; (C.7a.1) preserves it. Formula (C.7a.2) forces both winding
coordinates of the image of $g^2$ to vanish. If $u_0$ is its winding
coordinate at $z=0$, pushforward to $U$ gives $2u_0=2A$. Thus $A=0$.
We have proved zero pro-five abelianization for the full fixed group at $2$.

Next take the tangent $\overrightarrow{01}$, and let $g$ be a full fixed
element. The full cusp inertia $I_0\simeq\widehat\Z(1)$ is fixed by $H$.
Multiply $g$ by the inverse of the element of its pro-five Sylow subgroup
with the same $a$-coordinate. The resulting full fixed element $h$ has
$a$-coordinate zero. For every $n\ge1$, it therefore lies in the
geometric subgroup of
\[
 U_{5^n}=\mathbf P^1\setminus\bigl(\{0,\infty\}\cup\mu_{5^n}\bigr)
          \longrightarrow U,\qquad x=z^{5^n}.
\]
The chosen unit tangent at $z=0$ lifts the section over $H$. This follows
directly by substituting $x=z^{5^n}$ in the Puiseux fibre functor; all
$5^n$th roots of unity already lie in $F$. Denote the pro-five winding
coordinates of $h$ on this cover by $u_0$ and $u_\zeta$ for
$\zeta\in\mu_{5^n}$. Pushforward gives $5^nu_0=0$, so $u_0=0$, and
\[
 B(h)=\sum_{\zeta\in\mu_{5^n}}u_\zeta.                       \tag{C.7a.3}
\]
For any distinct $\zeta,\eta$ fill in $0$ and all other finite cusps.
The tangent then induces the ordinary point $0$ on
$\mathbf P^1\setminus\{\zeta,\eta,\infty\}$. In the coordinate
$t=(z-\zeta)/(\eta-\zeta)$, that point has
\[
 t=-\zeta/(\eta-\zeta),\qquad 1-t=\eta/(\eta-\zeta).
\]
Consequently both its Kummer classes on $H$ equal
$-\kappa(\eta-\zeta)$. This class is nonzero: the difference of distinct
five-power roots of unity has positive valuation, and (C.7a.1) applies
over $M=\Q_5(\mu_{5^n})$. Formula (C.7a.2) gives
$u_\zeta=u_\eta$. Thus (C.7a.3) makes $B(h)$ divisible by $5^n$ for
every $n$, and hence $B(h)=0$. For every $n$ all the $u_\zeta$ then vanish
as well.

It remains to record the stronger nonabelian quotient consequence.
Put $\Lambda=\Z_5[[\Z_5]]=\Z_5[[T-1]]$, where $T$ denotes conjugation
by $a$. The free pro-five presentation gives
\[
 N^{\rm ab}\simeq\Lambda\,b,\qquad Q\simeq\Lambda\rtimes\Z_5.
                                                               \tag{C.7a.4}
\]
For completeness, $N$ is the closed normal subgroup generated by $b$.
Its abelianization is therefore topologically generated by the conjugates
$a^i b a^{-i}$ and is a cyclic continuous $\Lambda$-module. The universal
property of the free pro-five group supplies the reverse map to
$\Lambda\rtimes\Z_5$ sending $a$ to the shift and $b$ to $1\in\Lambda$.
The resulting maps are inverse on the module generator and its translates,
which proves (C.7a.4). Under this identification, the coordinates
$(u_\zeta)$ at level $n$ are the coefficients of the image in
$\Lambda/(T^{5^n}-1)$. Their vanishing for every $n$ proves that the image
of $h$ in $\Lambda$ is zero: reduction modulo $5^r$ and $T^{5^n}-1$
for all $r,n$ defines the completed group-ring topology. Thus the tangent
fixed group has image contained in $\overline{\langle a\rangle}$ in $Q$;
cusp inertia gives the reverse inclusion.

We finish by strengthening the ordinary-point conclusion from $Q$ to
all of $P$. Use the Magnus homomorphism
\[
 P\longrightarrow 1+(X,Y)\Z_5\langle\!\langle X,Y\rangle\!\rangle,
 \qquad a\longmapsto1+X,\quad b\longmapsto1+Y.                  \tag{C.7a.5}
\]
The target has its $(5,X,Y)$-adic topology; it is pro-five, so the
homomorphism exists by freeness. It is injective. Indeed, for any finite
five-group quotient $D$ of $P$ and any $r\ge1$, the augmentation ideal
of $(\Z/5^r)[D]$ is nilpotent. To see this first modulo $5$, choose a
central element $z$ of order $5$. The quotient by $(z-1)$ is the group
algebra of the smaller five-group $D/\langle z\rangle$. Induction shows
that a power of the augmentation ideal is contained in $(z-1)$, whose
fifth power is zero; this proves nilpotence. Modulo $5^r$, a power of the
ideal is contained in $5(\Z/5^r)[D]$, so a further $r$th power vanishes.
Consequently the substitution $X\mapsto\bar a-1$, $Y\mapsto\bar b-1$
defines a continuous homomorphism from the series ring to $(\Z/5^r)[D]$.
It recovers the quotient map $P\to D$ on the Magnus image. An element
in the kernel of (C.7a.5) is therefore trivial in every finite five-group
quotient, hence trivial in $P$.

Choose $\sigma\in H$ for which $c=\kappa(2)(\sigma)\ne0$, using the
integral five-adic Kummer cocycle. On the series ring the tangential
action has
\[
 X\longmapsto X,\qquad Y\longmapsto Y+\text{terms of degree at least }3,
\]
because $f_\sigma\in[P,P]$. The path cocycle from the tangent to $2$
has expansion $1+cX+\text{terms of degree at least }2$: its $Y$
coefficient is $\kappa(-1)(\sigma)=0$. Hence the first nontrivial
homogeneous derivation of the point action is
\[
 \delta(X)=0,\qquad \delta(Y)=c(XY-YX),\qquad
 \delta(W)=c(XW-WX).                                          \tag{C.7a.6}
\]
Reversing a path convention reverses the common sign of this formula,
which does not alter its kernel.

Suppose the image of a full fixed element $g$ in $P$ were nontrivial,
and write its Magnus expansion as $1+G_n+G_{n+1}+\cdots$, with $G_n$
the first nonzero homogeneous polynomial. The quadratic-cover argument
already proved that the $a$-coordinate of $g$ is zero. Setting $Y=0$ in
its entire Magnus expansion therefore gives $1$; in particular $G_n$
has zero coefficient on the pure word $X^n$. Fixedness under $\sigma$,
in degree $n+1$, gives $c[X,G_n]=0$ by (C.7a.6), so $XG_n=G_nX$.
The centralizer of $X$ in the free associative polynomial algebra is
$\Z_5[X]$. One elementary verification is to compare coefficients of
$XG_n$ and $G_nX$: a word with nonzero coefficient must begin with $X$,
and moving that initial $X$ to its end preserves the coefficient.
Repeating would bring a $Y$ to the front if any $Y$ occurred, an
impossibility. Thus $G_n$ is a scalar multiple of $X^n$, contradicting
its zero pure-$X$ coefficient. This proves the ordinary-point assertion.

The argument controls images of fixed elements, not all invariants of
quotients. In fact the tangential $H$-action on $Q$ itself is trivial:
$a$ is fixed, and conjugation of $b$ by $f_\sigma\in[P,P]\subset N$
vanishes in $N^{\rm ab}$. Thus the tangent assertion gives a concrete
failure of surjectivity from full invariants to quotient invariants.
\end{proof}
